\subsection{Optimización específica para Grover mediante parametrización reducida}
\label{subsec:grover_reducido_impl}

Además del esquema general basado en bloques, se incorporó una variante especializada para las
familias de Grover consideradas en la evaluación. Esta variante aprovecha que, con
preparación uniforme y un conjunto fijo de estados marcados, la evolución ideal conserva dos
clases de amplitud: una común a los estados marcados y otra común a los estados no marcados.

En lugar de materializar el vector completo de tamaño $N$, la simulación mantiene únicamente
las amplitudes representativas de ambas clases, denotadas $u_t$ y $v_t$, junto con el valor
de $M$, número de estados marcados. La actualización de una iteración completa de Grover se
reduce entonces a la recurrencia
\[
\begin{bmatrix}
u_{t+1} \\
v_{t+1}
\end{bmatrix}
=
\begin{bmatrix}
\frac{N-2M}{N} & \frac{2(N-M)}{N} \\
-\frac{2M}{N} & \frac{N-2M}{N}
\end{bmatrix}
\begin{bmatrix}
u_t \\
v_t
\end{bmatrix},
\]
equivalente a una actualización de dimensión constante determinada únicamente por $N$ y $M$.
La derivación de esta expresión se detalla en el Apéndice~\ref{ap:grover_reducido}.

Esta optimización no reemplaza el esquema comprimido general ni constituye una alternativa general de
almacenamiento para otros circuitos. Su validez depende de que el algoritmo preserve la
simetría entre estados marcados y no marcados; por esa razón se utilizó como comparación
complementaria en Grover, donde permite estimar el beneficio adicional que puede obtenerse al
explotar directamente la estructura matemática del circuito.
